\documentclass{article}

% Document style
\usepackage{arxiv}
\usepackage{expl3}

% Encoding and fonts
\usepackage[utf8]{inputenc}   % UTF-8 support
\usepackage[T1]{fontenc}      % 8-bit T1 fonts

% Typography & formatting
\usepackage{microtype}        % Better spacing & kerning
\usepackage{nicefrac}         % Compact 1/2 symbols
\usepackage{booktabs}         % Professional-quality tables
\usepackage{caption}          % Better control of captions
\usepackage{float}            % Improved float handling

% Colors
\usepackage[dvipsnames]{xcolor} % Extended color set
\definecolor{gren}{HTML}{589283}

% Graphics & drawing
\usepackage{graphicx}
\usepackage{tikz}
\usetikzlibrary{shapes.multipart, positioning}

% Math packages
\usepackage{amsmath, amssymb, amsfonts, mathtools}

% Algorithms
\usepackage{algorithm}
\usepackage{algpseudocode}

% Boxes
\usepackage[most]{tcolorbox}

% Tables & lists
\usepackage{tabularx}
\usepackage{enumitem} % For compact/enhanced item lists

% References & citations
\usepackage[hidelinks]{hyperref}
\usepackage{url}
\usepackage{natbib} % has a nice set of citation styles and commands
    \bibliographystyle{plainn}
    \renewcommand{\bibsection}{\subsubsection*{References}}
\usepackage{doi}

% Page layout
\usepackage{geometry}
\geometry{margin=1in}

% Notes (optional, can remove if not needed)
\usepackage[colorinlistoftodos]{todonotes}
\setlength{\marginparwidth}{2cm}

% Dummy text (remove in final draft)
\usepackage{lipsum}
\newcolumntype{M}[1]{>{\raggedright\arraybackslash}m{#1}}


% Light Blue Rounded Box
\newtcolorbox{bluebox}[1][]{
    colback=blue!2,       % light blue background
    colframe=blue!50,      % darker blue border
    rounded corners,
    boxrule=0.8pt,
    arc=4mm,
    title=#1
}




\title{Pre-training in Reinforcement Learning: Data, Curiosity and Structure}

\author{ Jasmeet Kaur \\ 
}

% Uncomment to remove the date
\date{}

\begin{document}
\maketitle

% citation check
% \section*{Conference Submission Timeline}

% \rowcolors{2}{gray!10}{white}
% \begin{tabular}{|l|l|c|}
% \hline
% \textbf{Date / Week} & \textbf{Task} & \textbf{Status} \\
% \hline-
% July 16   & Literature Review A & \cellcolor{green!40} Done \\
% Aug 8  & Literature Review B & \cellcolor{green!40} Done \\
% Sept 8 & Literature Review C & \cellcolor{green!40} Done  \\
% Sept 15 & Write Experiments, Conclusion & \cellcolor{gray!40} To Do \\
% Sept 20 & Literature Review c & \cellcolor{gray!50} To Do  \\
% Sept 22  & Final Polish & \cellcolor{gray!40} To Do \\
% Sept 25    & Submission deadline & \cellcolor{gray!40} To Do  \\
% \hline
% \end{tabular}



\section{Introduction}

Label-free learning has enabled recent progress in learning from large amounts of data. Advancements in Large Language Models showcase the effectiveness of learning the generative process from huge amounts of text data. This also pushed learning without labels and learning from auxiliary labels in some of the sub-fields of Reinforcement Learning (\cite{park2024foundation}, \cite{rudolph2024learning}, \cite{li2023accelerating}, \cite{ma2022vip}, etc.). This article provides a focused study of different works unified by data-driven objectives. It filters the entire discussion of using data to learn representations into two perspectives. First is learning driven by curiosity to explore the dataset, and the second focuses on learning the underlying structure. 

This article covers works that focus on learning from data that is invariant to task-relevant information (\cite{zisselman2023explore}, \cite{wilcoxson2024leveraging}, \cite{rudolph2024learning}, \cite{laskin2020curl}, etc.). The discussion highlights the type of underlying structure in the data, what are the ways to learn about this structure, and how to adapt these learned insights to solve downstream tasks with minimal intervention. The aim is to give an integrated understanding of these works and where this may lead us in the future.

Prevailing practices demonstrate that the availability of data has made its use and scientific consideration inevitable. Although not all data is good, a lot of data used for training RL agents comes from unknown tasks, suboptimal, and incomplete interactions (\cite{hansen2022bisimulation}, \cite{park2024foundation}, \cite{li2023accelerating}). Methods to pre-train using this data are attractive because solving new tasks from scratch is not efficient. Utilizing data is beneficial in cases where interacting with the environment during deployment is not possible. 

Most of the papers covered under the discussion presented here are works related to unsupervised and self-supervised pre-training for learning representations in RL. Some works that lay the foundation of such approaches have also been described in the relevant sections. The entire analysis has been divided into four main sections. Each section covers a pre-training objective that is motivated by either of the two main factors described above. This article does not propose a novel method to solve any problem, but offers an organized discussion.
 
The focus for each section is as follows: exploring datasets for solving downstream tasks with minimal or no adaptation, learning representations that capture structure in environmental dynamics, learning representations that capture similar functional behavior within an environment, and learning representations that capture structure between similar and dissimilar behaviors. The objectives include maximizing entropy, finding temporal structure, learning successor features, learning with contrastive objectives, and bisimulation metrics. These methods provide overlapping answers to some important research questions. These include finding the right type of representations to learn from data and understanding the trade-offs of different pre-training objectives for efficient downstream adaptation. 

\section{Problem Statement}

All problem settings presented operate in an MDP (Markov-decision process (\cite{bellman1958dynamic}])) framework or some extension of this framework. The environment is described as an MDP $M$ with state space $\mathcal{S}$, action space $\mathcal{A}$, transition dynamics $P$, and discount factor $\gamma$. A policy $\pi(a|s)$ describes what action to take given a state $s$. Each task is defined via a reward function $r$ that maps state-action pairs to real-valued rewards. In some cases, an offline dataset of trajectories is available. This dataset, denoted as $\mathcal{D}$, consists of $N$ trajectories, where each trajectory $\tau_i$ is a sequence of tuples $(s_t^i, a_t^i, s_{t+1}^i)$ for $t = 0$ to $T_i$, with $T_i$ being the horizon length of the trajectory. These trajectories can be collected using any policy, whether optimal or suboptimal, and may involve unknown reward functions.


\section{Methods}

This section focuses on four different approaches for representation learning in RL. These representations aim at few-shot adaptations to downstream tasks. The section describes the different pre-training objectives, adaptation or fine-tuning strategies for new tasks, implementation choices, and open questions. The goal is to cover the strengths and impact of these contributions.

\subsection{Exploration-based Representations}
Incorporating exploration as an objective can be used to learn diverse skills from data. Exploration-based works focus on invariance to task rewards. This can be done for either the offline data available or for gathering new data through online interactions. 
\cite{zisselman2023explore} argues that exploring a domain is harder than training a policy to maximize reward in an environment. A pre-training exploratory policy can be used for zero-shot RL (\cite{zisselman2023explore}). This can be done by training an exploratory policy that maximizes the entropy of a marginal state visitation distribution for a finite horizon $T$. This distribution is defined as:
\begin{equation}
d_{T,\pi}(s) = \frac{1}{T} \sum_{t=0}^{T} d_{t,\pi}(s),
\end{equation}
where $d_{t,\pi}(s)$ denotes the $t$-step state distribution when following policy $\pi$. The entropy objective helps move towards the states that the policy has not encountered in the episode. \textbf{\textit{D}} The objective maximized over a horizon T, becomes (\cite{zisselman2023explore}):
\begin{align}
R_T(\pi) 
&= \mathbb{E}_{\mathcal{M} \sim \hat{P}(\mathcal{M})}
  \left[ \emph{H}( d_{T, \pi}(s)) \right] \\
&= \mathbb{E}_{\mathcal{M} \sim \hat{P}(\mathcal{M})}
  \left[
    \emph{H}\!\left( \frac{1}{T} \sum_{t=0}^{T} d_{t,\pi}(s) \right)
  \right].
\end{align}

The expectation is over different MDPs sampled from a distribution. This objective learns a policy that equally visits all the states during the episodes (\cite{zisselman2023explore}). Entropy can be computed by a particle-based k-nearest neighbor approximation (\cite{zisselman2023explore}, \cite{yarats2021reinforcement}). For practical implementation, \cite{zisselman2023explore} treats the objective as a reward to be optimized for by training an agent with PPO. But learning an exploratory policy is not enough, and disentangling exploration from underlying dynamics is important. One way to do this is to train an ensemble of policies and use the exploratory policy to find actions in the environment when the ensemble fails to agree on an action (\cite{zisselman2023explore}).

Evidently, the entropy maximizing policy shows a small generalization gap on the unknown tasks for effective zero-shot generalization. In Zero-shot generalization, a policy is trained on a set of tasks and tested on unknown tasks. The generalization quality of this policy is tested by training agents on 200, 500, and 1000 of Procgen's Maze, Jumper, and Miner (\cite{zisselman2023explore}) and such exploration learns reward invariance. 

A sample-efficient approach with a similar motivation can be used for representation learning. Learn a visual representation without task rewards in a self-supervised way. The pre-training phase has two parts. First, learn a set of prototypical embeddings that form the fundamental basis of a low-dimensional latent space. Second, use these prototypes to maximize an entropy-based intrinsic reward, encouraging exploration of the environment. The intrinsic reward is similar to the one in \cite{zisselman2023explore}, but instead of states, it uses encoded observations. \cite{yarats2021reinforcement} shows the performance of such an approach on downstream tasks. It appends the intrinsic reward with the task reward to train a SAC agent for DeepMind Control. 

Both methods assume interaction with the environment in the pre-training phase, with no fine-tuning for downstream adaptations. But diverse behavior from offline data can also be learnt in a task-agnostic way. The idea is to optimize for diverse temporal behavior in the pre-training phase. Learn a state representation $\phi$ that captures the temporal structure from unlabeled trajectories by using the equivalence between temporal distance and the optimal goal-conditioned value function. \cite{park2024foundation} uses this representation to learn a policy that spans the latent space and captures diverse skills in the offline data. Few-shot adaptation to the downstream can be achieved by learning the task-dependent information.

Curiosity-driven objectives can be used to improve online data collection and online exploration for new tasks. Such objectives are effective in solving sparse reward tasks at test time (\cite{burda2018exploration}, \cite{li2023accelerating}). The unlabeled data is labeled with optimistic rewards using RND (\cite{burda2018exploration}) to guide exploration to unknown states. An RL agent trained with this data can be used to collect data through online interactions and update the reward model using this online experience. This offers an iterative approach for fine-tuning the agent with updated rewards. \cite{li2023accelerating} evaluates this method for rapid exploration in challenging sparse-reward domains such as the AntMaze domain, Adroit hand manipulation, and a visually simulated robotic manipulation domain. 

Offline data can be used with such an objective to improve online exploration. The key idea of this is to extract the low-level skills from the offline dataset using a variational encoder. A learnt prior conditions the posterior to stay close to the offline dataset. Learn a low-level policy to take actions around the dataset trajectories and learn a high-level policy through online interactions. This method is effective for tasks that need exploration to be solved (\cite{wilcoxson2024leveraging}). Both these approaches improve online exploration by using unlabeled data and implicitly capturing invariance.

\subsection{Structured Representations}
A few approaches present a way to learn representations that explicitly exploit structure in the environment. The structure across different tasks remains the same when the environment dynamics are the same. This section focuses on one way to learn such invariance. 

\subsubsection{Successor Features}
Successor Features were used to do transfer in RL by exploiting the underlying structure in dynamics and performing well on new tasks (\cite{barreto2017successor}). If the expected one-step reward associated with transition $(s, a, s')$ is decomposed into features represented by $\phi$ and weights $w$, such a representation can be used to compute the $Q$-value function. This $Q$-value function can be decomposed into the task-relevant information and cumulant $\psi$, which is called Successor Features. Once Successor Features are learnt for different tasks, Generalized Policy Iteration is used to learn a policy that does as good as any prior policy on this new task (\cite{barreto2017successor}). 

A key focus of an earlier application of this idea is that with SF and GPI, no task-specific information is shared to learn policies (\cite{borsa2018universal}). SF$\&$GPI can be combined with the UVFAs (\cite{schaul2015universal}) to learn the value functions over the task encodings and learn a span over the policies. While this offers good generalization, the performance of the policies on a new task depends on how the policies are sampled. An alternative approach for learning shared knowledge across different tasks is using an attention head to learn successor features (\cite{carvalho2023composing}). A more robust way to generate approximations of successor features is to learn a distribution of the returns (\cite{carvalho2023combining}). Finally, such structural decomposition can be used to learn to adapt to tasks with non-Markovian rewards (\cite{moskovitz2021first}). The successor representations capture the expected steps to reach a state from the current state for the first time. 


\subsubsection{Successor Measure}
The successor feature decomposition can be extended to state-visitation distributions. The idea is to decompose the Q-function into the successor measure and the reward. The successor measure encodes the probability of going to the next state under the policy (\cite{touati2021learning}). The entire family of policies can be parametrized with z, and the Q-function is used to learn the optimal policies. As in \cite{touati2021learning}, learning the successor measure representation for all z allows for adaptation to any reward function. For each $z$, find $F_z$ and $B$ such that the successor measure $M^{\pi}_z$ is $M^{\pi}_z = F_z^{\top} B$. Then, for a reward function $r$, the action-value function is given by $Q^{\pi}_{z,r} = F_z^{\top} B\,r$ (\cite{touati2021learning}).

Learning the span of policies does not have to be restricted to optimal policies (\cite{agarwal2024proto}). Successor Measure can be used to learn representations in the entire space of policies. \cite{agarwal2024proto} transforms the constrained optimization problem following from the Bellman Flow equations of successor measure into a single-player game.

Successor Features and Successor Measure offer a concise way to represent the underlying structure explicitly. Their ability to transfer to downstream tasks relies on how the space of policies is encoded with this information, making it an important area of closer examination. 


\subsection{Bisimulation-based Representations}
Bisimulation metrics capture invariance to task-irrelevant features. They group functionally similar states in a latent space, and capturing this similarity in an MDP can be useful for generalization over large state spaces (\cite{ferns2004metrics}). In theory, Bisimulation defines an equivalence relation over the state-space. Partitioning the state-space under this relation is difficult for practical application (\cite{ferns2004metrics}). To deal with this, a semi-metric over this relation is defined to capture how similar two states are.

While bisimulation metrics are used to learn representations from image observations in reinforcement learning now, putting bisimulation to practical use has come a long way. Issues related to its efficient computation and finding upper bounds for arbitrary policies proved worthwhile. A sampling-based approximation to the joint distribution of the transition dynamics made it easier to implement for learning approximation to a large-state space (\cite{castro2021mico}). A common problem that emerges is that of representation collapse, where two dissimilar states end up with the same representations. This is more significant in sparse reward settings. One way to mitigate this is to use cosine distance to measure the distance between two latent states (\cite{zang2022simsr}). 
\begin{align}
\cos_{\phi}(x, y) 
&= 1 - \frac{\phi(s)^\top \phi(t)}
        {\|\phi(s)\| \cdot \|\phi(t)\|},
\label{eq:cosine} \\[1ex]
F^{\pi}_{\cos\phi}(s, t) 
&= \big| r^{\pi}_s - r^{\pi}_t \big|
   + \gamma \, \mathbb{E}_{s' \sim \hat{P}^{\pi}_s, \; t' \sim \hat{P}^{\pi}_t}
      \big[ \cos_{\phi}(s', t') \big].
\label{eq:simsr}
\end{align}

Aforementioned methods provide objectives to learn effective representations efficiently using bisimulation metrics. The following subsections describe three domains where bisimulation metrics have been adapted for specific problems. More important to this discussion are works on learning action representations and learning from offline data with no action information. In general, these works show the effectiveness of capturing this behavioral similarity in different settings. These adaptations have been successful in solving complex tasks such as robot manipulation (\cite{hansen2022bisimulation}, \cite{shi2024self}). 

\subsubsection{Goal-conditioned Bisimulation Relations}
Bisimulation can result in the effective transfer of skills across analogous tasks (\cite{hansen2022bisimulation}). Given an MDP $M$, \emph{goal-conditioned bisimulation relation}  $\mathcal{B}$ can be defined (\cite{hansen2022bisimulation}). For all state - goal pairs $(s, g_s), (t, g_t) \in S \times G$ that are equivalent under 
$\mathcal{B}$, the following conditions hold:
\begin{align}
    R(s, a, g_s) &= R(t, a, g_t), \quad \forall a \in A \label{eq:reward-cond} \\
    P(G \mid s, a) &= P(G \mid t, a), \quad \forall a \in A, \; G \subseteq S. \label{eq:transition-cond}
\end{align}

For practical implementation, an \emph{on-policy} version of this relation gives rise to a paired-state metric:
\begin{align}
    d\big((s,g_s), (t,g_t)\big) 
    &= \big| R(s, \pi(s,g_s), g_s) - R(t, \pi(t,g_t), g_t) \big| \notag \\
    &\quad + W_2 \Big( P(\cdot \mid s, \pi(s,g_s)), \;
    P(\cdot \mid t, \pi(t,g_t)) \Big), \label{eq:metric}
\end{align}
where $\pi$ is the goal-conditioned policy. This objective can be used to learn state-goal representations and train an offline goal-conditioned policy. \cite{hansen2022bisimulation} used a dataset collected using a noisy expert to show skill transfer in manipulation tasks. 
\begin{align}
L_{\varphi} &= \biggl( \bigl\|
    \phi(s, g_s) - \phi(t, g_t)
\bigr\|_1
- \bigl\| R_s - R_t \bigr\|_2
- \gamma \bigl\|
    \overline{\phi}(s', g_s') - \overline{\phi}(t', g_t')
\bigr\|_2
\biggr)^2
\end{align}

A bisimulation objective can be appended with a forward model error-based intrinsic reward to improve exploration. This enhances exploration in the latent space for learning a bisimulation-based representation and has been empirically shown to learn robust state representations for downstream goal-conditioned tasks (\cite{kemertas2021towards}). 


\subsubsection{Learning from Offline Data}
Bisimulation-based representations can be learnt using an offline dataset. Learning from an offline dataset is notoriously difficult using function approximation. The representations learnt using the bisimulation metric can result in a value function that does not diverge. These representations stabilize TD-learning and are Bellman complete (\cite{pavse2024stable}). \cite{pavse2024stable} defines a $\pi$-bisimilarity kernel $k^{\pi_e} : S \times S \to \mathbb{R}$ for pairs of state-actions and under $\pi_e$:
\begin{align}
k^{\pi_e}(s,a_s; t,a_t) 
&= k_1(s,a_s; t,a_t) 
   + \gamma \, k_2\big(k^{\pi_e}\big)\big(P^{\pi_e}(\cdot|s,a_s), P^{\pi_e}(\cdot|t,a_t)\big),
\end{align}
Here, $k_1$ measures short-term similarity based on the rewards received. $k_2$ measures long-term similarity between probability distributions by evaluating similarity between samples of the distributions according to $k^{\pi_e}$.

To use the bisimulation metric for offline datasets, it must learn from an incomplete dataset with missing transitions. While Implicit Q learning can be used to mitigate this problem (\cite{hansen2022bisimulation}), the error in the estimated value of bisimulation operator can be reduced by direct use of an expectile operator (\cite{zang2023understanding}).
\begin{align}
F^{\pi_\beta} \phi^{\pi_\beta}(s, t) 
&:= \arg\min_{\phi^{\pi_\beta}} 
\;\; \mathbb{E}_{a_s \sim \pi_\beta(\cdot|s), \; a_t \sim \pi_\beta(\cdot|t)} \bigg[
    \tau \, [\hat{\varepsilon}]_+^2 
    + (1 - \tau) \, [-\hat{\varepsilon}]_+^2
\bigg], \label{eq:expectile_operator} \\
\hat{\varepsilon} 
&= \mathbb{E}_{s' \sim T^{\pi}(s), \; t' \sim T^{\pi}(t)} 
\Big[ \; \big| r(s,a_s) - r(t,a_t) \big| 
    + \gamma \, \overline{\phi}^{\pi_\beta}(s', t') 
    - \phi^{\pi_\beta}(s,t) \;\Big].
\label{eq:residual}
\end{align}

Finally, for real-world deployment of bisimulation-based representations, naively training with adversarial states and actions does not transfer well to the downstream tasks (\cite{yin2024representation}). Learn robust representation with perturbed states and goals using a contrastive objective. 

\subsubsection{Action Invariance} 
Bisimulation-based metrics enable capturing invariance in behavior by learning action representations. Such long-term action representations can be learnt in a self-supervised way using bisimulation (\cite{shi2024self}). The state-conditional action chunk (denoted by c) bisimulation metric is a function 
$d: \mathcal{S} \times \mathcal{C} \times \mathcal{C} \to \mathbb{R}_{\ge 0}$ such that
\begin{align}
d(c_i, c_j \mid s_t) 
&= R^{c_i}_{s_t} - R^{c_j}_{s_t} 
   + \gamma \, W_2 \big( P^{c_i}_{s_t}, P^{c_j}_{s_t}; d_{c} \big),
\label{eq:bisim_metric}
\end{align}
where $d$ is a pseudometric
and $W_2$ is the 2nd Wasserstein distance between two distributions.  
Here, $R^c_{s_t}$ represents the cumulative discounted reward for executing chunk $c$ starting at $s_t$ and 
$P^c_{s_t}$ represents the distribution of $s_{t+k}$ after executing $c$ from $s_t$.
Such representations are effective for solving complex tasks and this is experimentally shown with 7DOF ARM Control (\cite{shi2024self}).

Learning action representations is not restricted to online interactions but can also be achieved from an offline dataset. To stabilize learning close to the behavioral policy and mitigate distribution mismatch, the following objective can be used (\cite{gu2022learning}):
\begin{align}
L(\phi) 
&= \mathbb{E}_{s, a_i, r \sim \mathcal{D}, \; a_j \sim \mathcal{A}} 
   \Big[ \big\| \phi(a_{i}) - \phi(a_{j}) \big\|_1 - \hat{d}(a_i, a_j \mid s) \Big]^2,
\label{eq:mod_obj}
\end{align}
where
\begin{align}
\hat{d}(a_i, a_j \mid s) 
&= \big| r_{i} - \hat{R}(s, \phi(a_{j})) \big|
   + \gamma \, W_2 \big( \hat{P}(\cdot \mid s,  \phi(a_{i})), \hat{P}(\cdot \mid s,  \phi(a_{i})) \big)
   + p \cdot \hat{I}_\beta(a_j \mid s).
\label{eq:estimated_d}
\end{align}
Here, $d(a_i, a_j \mid s)$ is an estimate of the true distance between actions conditioned on the same state $s$. $\hat{R}$ and $\hat{P}$ are the learned reward and transition models, trained separately. $\hat{I}_\beta(a_j \mid s)$ is a trainable model predicting whether $a_j$ is out-of-distribution.

While the method above explicitly handles the behavioral distribution, a similar objective for single-step action representation, denoted by $\phi$, can be learnt. Such an action-bisimulation metric is defined as (\cite{rudolph2024learning}):
\begin{align}
d(s_i, s_j, \phi, \varphi) 
&= (1-c) \cdot \| \phi_(s_i) - \phi_(s_j) \|_1 
   + c \cdot \mathbb{E}_{a \sim \mathcal{U}(\mathcal{A})} 
       \Big[ W_1 \big( f(\varphi(s_i), a), f(\varphi_(s_j), a) \big) \Big],
\label{eq:da_bisim}
\end{align}
where $c \in [0,1]$ balances the contributions of the state and action-conditional terms, $W_1$ denotes the 1st Wasserstein distance and $f$ is the trained  forward model. The state representations are learnt by minimizing the $L_1$ distance between embedded representations 
$\varphi_(s_i)$ and $\varphi_(s_j)$ to the action-bisim metric (\cite{rudolph2024learning}):
\begin{align}
L(\mathcal{D}) 
&= \frac{1}{N} \sum_{s_i, s_j \sim \mathcal{D}} 
   \Big| \| \varphi_(s_i) - \varphi_(s_j) \|_1 
   - d_{\text{a-bisim}}(s_i, s_j, \psi, \varphi) \Big|.
\label{eq:da_bisim_loss}
\end{align}

Bisimulation effectively captures the similarity between states. It can be adapted for offline datasets, online exploration, and learning action representations. 

\subsection{Contrastive Learning based Representations}

Contrastive Learning objectives have been used for unsupervised RL. It is a framework to learn representations by exploiting the structure between similar and dissimilar pairs of input. This can be achieved in an unsupervised way by performing a dictionary lookup task wherein the positives and negatives represent a set of keys with respect to a query. There are many ways in which such an objective can be captured, and one of the most adapted ones is the InfoNCE loss function. 

\begin{align}
\mathcal{L}_{\text{contrastive}} &= 
- \log \frac{\exp(\text{sim}(q, k^+)/\tau)}
       {\sum_{k \in \mathcal{K}} \exp(\text{sim}(q, k)/\tau)}
\end{align}

Here, $q$ is the query embedding, $k^+$ is the positive key embedding, $\mathcal{K}$ is the set of all keys in the mini-batch (positives and negatives), $\text{sim}(x, y) = \frac{x \cdot y}{\|x\| \|y\|}$ denotes the cosine similarity, and $\tau$ is the temperature hyperparameter.

Contrastive Learning can be used to train an RL agent over representations learnt from image-based observations (\cite{laskin2020curl}). It can be used to improve the robustness of bisimulation representation for transfer to downstream tasks by using it to perturb the negative samples (\cite{yin2024representation}). $P_{trb}$ in the following objective denotes the learnt perturbation parametrized by $\theta$ (\cite{yin2024representation}):
\begin{align}
\mathcal{L}(\theta) &= - \phi \big( P_{trb}^i(s), P_{trb}^i(g) \big)^\top \phi\big( \langle s, g \rangle\big)^{-}
\end{align}

Finally, Contrastive Learning is effective for representations over offline image data (\cite{ma2022vip}). Learn representation over states using a goal-conditioned offline pre-training objective. The objective minimizes the distance between the goal-conditioned state-occupancy distribution of the policy and the data distribution. The dual of this objective yields a contrastive RL objective. The objective is:
\begin{align}
\max_{\pi_T, \phi} \; & \mathbb{E}_{\pi_T} \Bigg[ \sum_t \gamma^t r(o; g) \Bigg] \nonumber \\
& - D_{\mathrm{KL}}\big(d_{\pi_T}(o, a_T; g) \,\|\, d_D(o, \tilde{a}_T; g)\big),
\end{align}
where $d_{\pi_H}(o, a_H; g)$ is the distribution over observations and actions visited by the policy $\pi_H$ conditioned on goal $g$. $d_D(o, \tilde{a}_H; g)$ is the distribution over observations and dummy actions $\tilde{a}_H$ in the dataset $D$, conditioned on goal $g$. This method has been effective for zero-shot generalization in goal-conditioned reinforcement learning (\cite{ma2022vip}).

\section{Conclusion}

This article provides an overview of recent research trends in RL. It looks at different works using unsupervised and self-supervised methods of learning in RL and provides a unification of the core objectives. It covers approaches based on exploration, successor features, bisimulation, and contrastive learning, highlighting the importance of curiosity and structure in RL. This synthesis of reviewed work indicates certain key insights. Curiosity-based approaches capture a notion of diversity and learns to disentangle this diversity for downstream tasks. In general, how exploration is related to representation learning for downstream performance deserves additional study. On the other other, Successor Features and Successor Measures capture structure explicitly. Structure in the environment evidently dictates the structure in the solution space, but this interdependence needs to be looked into further. Bisimulation learns behavioral similarity, but computing the metric using probability distributions is a bottleneck. Contrastive Learning provides a sample-efficient way to learn representations that are close to positive samples and distant from negative samples. 



\section{Use of Generative AI}
LaTeX code for equations has been modified from templates generated with the help of ChatGPT (OpenAI). In addition, it has been used to assist in finding word phrasing. The text was proofread using Grammarly’s basic grammar and style suggestions, without the use of AI-generated content. 





\begin{thebibliography}{27}
\providecommand{\natexlab}[1]{#1}
\providecommand{\url}[1]{\texttt{#1}}
\expandafter\ifx\csname urlstyle\endcsname\relax
  \providecommand{\doi}[1]{doi: #1}\else
  \providecommand{\doi}{doi: \begingroup \urlstyle{rm}\Url}\fi

\bibitem[Agarwal et~al.(2024)Agarwal, Sikchi, Stone, and Zhang]{agarwal2024proto}
Siddhant Agarwal, Harshit Sikchi, Peter Stone, and Amy Zhang.
\newblock Proto successor measure: Representing the behavior space of an rl agent.
\newblock \emph{arXiv preprint arXiv:2411.19418}, 2024.

\bibitem[Barreto et~al.(2017)Barreto, Dabney, Munos, Hunt, Schaul, van Hasselt, and Silver]{barreto2017successor}
Andr{\'e} Barreto, Will Dabney, R{\'e}mi Munos, Jonathan~J Hunt, Tom Schaul, Hado~P van Hasselt, and David Silver.
\newblock Successor features for transfer in reinforcement learning.
\newblock \emph{Advances in neural information processing systems}, 30, 2017.

\bibitem[Borsa et~al.(2018)Borsa, Barreto, Quan, Mankowitz, Munos, Van~Hasselt, Silver, and Schaul]{borsa2018universal}
Diana Borsa, Andr{\'e} Barreto, John Quan, Daniel Mankowitz, R{\'e}mi Munos, Hado Van~Hasselt, David Silver, and Tom Schaul.
\newblock Universal successor features approximators.
\newblock \emph{arXiv preprint arXiv:1812.07626}, 2018.

\bibitem[Burda et~al.(2018)Burda, Edwards, Storkey, and Klimov]{burda2018exploration}
Yuri Burda, Harrison Edwards, Amos Storkey, and Oleg Klimov.
\newblock Exploration by random network distillation.
\newblock \emph{arXiv preprint arXiv:1810.12894}, 2018.

\bibitem[Carvalho et~al.(2023{\natexlab{a}})Carvalho, Filos, Lewis, Singh, et~al.]{carvalho2023composing}
Wilka Carvalho, Angelos Filos, Richard~L Lewis, Satinder Singh, et~al.
\newblock Composing task knowledge with modular successor feature approximators.
\newblock \emph{arXiv preprint arXiv:2301.12305}, 2023{\natexlab{a}}.

\bibitem[Carvalho et~al.(2023{\natexlab{b}})Carvalho, Saraiva, Filos, Lampinen, Matthey, Lewis, Lee, Singh, Jimenez~Rezende, and Zoran]{carvalho2023combining}
Wilka~Carvalho Carvalho, Andre Saraiva, Angelos Filos, Andrew Lampinen, Loic Matthey, Richard~L Lewis, Honglak Lee, Satinder Singh, Danilo Jimenez~Rezende, and Daniel Zoran.
\newblock Combining behaviors with the successor features keyboard.
\newblock \emph{Advances in neural information processing systems}, 36:\penalty0 9956--9983, 2023{\natexlab{b}}.

\bibitem[Castro et~al.(2021)Castro, Kastner, Panangaden, and Rowland]{castro2021mico}
Pablo~Samuel Castro, Tyler Kastner, Prakash Panangaden, and Mark Rowland.
\newblock Mico: Improved representations via sampling-based state similarity for markov decision processes.
\newblock \emph{Advances in Neural Information Processing Systems}, 34:\penalty0 30113--30126, 2021.

\bibitem[Ferns et~al.(2004)Ferns, Panangaden, and Precup]{ferns2004metrics}
Norm Ferns, Prakash Panangaden, and Doina Precup.
\newblock Metrics for finite markov decision processes.
\newblock In \emph{UAI}, volume~4, pages 162--169, 2004.

\bibitem[Gu et~al.(2022)Gu, Zhao, Chen, Li, Hao, and An]{gu2022learning}
Pengjie Gu, Mengchen Zhao, Chen Chen, Dong Li, Jianye Hao, and Bo~An.
\newblock Learning pseudometric-based action representations for offline reinforcement learning.
\newblock 2022.

\bibitem[Hansen-Estruch et~al.(2022)Hansen-Estruch, Zhang, Nair, Yin, and Levine]{hansen2022bisimulation}
Philippe Hansen-Estruch, Amy Zhang, Ashvin Nair, Patrick Yin, and Sergey Levine.
\newblock Bisimulation makes analogies in goal-conditioned reinforcement learning.
\newblock In \emph{International Conference on Machine Learning}, pages 8407--8426. PMLR, 2022.

\bibitem[Kemertas and Aumentado-Armstrong(2021)]{kemertas2021towards}
Mete Kemertas and Tristan Aumentado-Armstrong.
\newblock Towards robust bisimulation metric learning.
\newblock \emph{Advances in Neural Information Processing Systems}, 34:\penalty0 4764--4777, 2021.

\bibitem[Laskin et~al.(2020)Laskin, Srinivas, and Abbeel]{laskin2020curl}
Michael Laskin, Aravind Srinivas, and Pieter Abbeel.
\newblock Curl: Contrastive unsupervised representations for reinforcement learning.
\newblock In \emph{International conference on machine learning}, pages 5639--5650. PMLR, 2020.

\bibitem[Li et~al.(2023)Li, Zhang, Ghosh, Zhang, and Levine]{li2023accelerating}
Qiyang Li, Jason Zhang, Dibya Ghosh, Amy Zhang, and Sergey Levine.
\newblock Accelerating exploration with unlabeled prior data.
\newblock \emph{Advances in Neural Information Processing Systems}, 36:\penalty0 67434--67458, 2023.

\bibitem[Ma et~al.(2022)Ma, Sodhani, Jayaraman, Bastani, Kumar, and Zhang]{ma2022vip}
Yecheng~Jason Ma, Shagun Sodhani, Dinesh Jayaraman, Osbert Bastani, Vikash Kumar, and Amy Zhang.
\newblock Vip: Towards universal visual reward and representation via value-implicit pre-training.
\newblock \emph{arXiv preprint arXiv:2210.00030}, 2022.

\bibitem[Moskovitz et~al.(2021)Moskovitz, Wilson, and Sahani]{moskovitz2021first}
Ted Moskovitz, Spencer~R Wilson, and Maneesh Sahani.
\newblock A first-occupancy representation for reinforcement learning.
\newblock \emph{arXiv preprint arXiv:2109.13863}, 2021.

\bibitem[Park et~al.(2024)Park, Kreiman, and Levine]{park2024foundation}
Seohong Park, Tobias Kreiman, and Sergey Levine.
\newblock Foundation policies with hilbert representations.
\newblock \emph{arXiv preprint arXiv:2402.15567}, 2024.

\bibitem[Pavse et~al.(2024)Pavse, Chen, Xie, and Hanna]{pavse2024stable}
Brahma~S Pavse, Yudong Chen, Qiaomin Xie, and Josiah~P Hanna.
\newblock Stable offline value function learning with bisimulation-based representations.
\newblock \emph{arXiv preprint arXiv:2410.01643}, 2024.

\bibitem[Rudolph et~al.(2024)Rudolph, Chuck, Black, Lvovsky, Niekum, and Zhang]{rudolph2024learning}
Max Rudolph, Caleb Chuck, Kevin Black, Misha Lvovsky, Scott Niekum, and Amy Zhang.
\newblock Learning action-based representations using invariance.
\newblock \emph{arXiv preprint arXiv:2403.16369}, 2024.

\bibitem[Schaul et~al.(2015)Schaul, Horgan, Gregor, and Silver]{schaul2015universal}
Tom Schaul, Daniel Horgan, Karol Gregor, and David Silver.
\newblock Universal value function approximators.
\newblock In \emph{International conference on machine learning}, pages 1312--1320. PMLR, 2015.

\bibitem[Shi et~al.(2024)Shi, Jianye, Tang, Dong, and Zheng]{shi2024self}
Lei Shi, HAO Jianye, Hongyao Tang, Zibin Dong, and Yan Zheng.
\newblock Self-supervised bisimulation action chunk representation for efficient rl.
\newblock In \emph{Neurips Safe Generative AI Workshop 2024}, 2024.

\bibitem[Touati and Ollivier(2021)]{touati2021learning}
Ahmed Touati and Yann Ollivier.
\newblock Learning one representation to optimize all rewards.
\newblock \emph{Advances in Neural Information Processing Systems}, 34:\penalty0 13--23, 2021.

\bibitem[Wilcoxson et~al.(2024)Wilcoxson, Li, Frans, and Levine]{wilcoxson2024leveraging}
Max Wilcoxson, Qiyang Li, Kevin Frans, and Sergey Levine.
\newblock Leveraging skills from unlabeled prior data for efficient online exploration.
\newblock \emph{arXiv preprint arXiv:2410.18076}, 2024.

\bibitem[Yarats et~al.(2021)Yarats, Fergus, Lazaric, and Pinto]{yarats2021reinforcement}
Denis Yarats, Rob Fergus, Alessandro Lazaric, and Lerrel Pinto.
\newblock Reinforcement learning with prototypical representations.
\newblock In \emph{International Conference on Machine Learning}, pages 11920--11931. PMLR, 2021.

\bibitem[Yin et~al.(2024)Yin, Wu, Liu, Fang, Zhao, Huang, and Ruan]{yin2024representation}
Xiangyu Yin, Sihao Wu, Jiaxu Liu, Meng Fang, Xingyu Zhao, Xiaowei Huang, and Wenjie Ruan.
\newblock Representation-based robustness in goal-conditioned reinforcement learning.
\newblock In \emph{Proceedings of the AAAI Conference on Artificial Intelligence}, volume~38, pages 21761--21769, 2024.

\bibitem[Zang et~al.(2022)Zang, Li, and Wang]{zang2022simsr}
Hongyu Zang, Xin Li, and Mingzhong Wang.
\newblock Simsr: Simple distance-based state representations for deep reinforcement learning.
\newblock In \emph{Proceedings of the AAAI conference on artificial intelligence}, volume~36, pages 8997--9005, 2022.

\bibitem[Zang et~al.(2023)Zang, Li, Zhang, Liu, Sun, Islam, Tachet~des Combes, and Laroche]{zang2023understanding}
Hongyu Zang, Xin Li, Leiji Zhang, Yang Liu, Baigui Sun, Riashat Islam, Remi Tachet~des Combes, and Romain Laroche.
\newblock Understanding and addressing the pitfalls of bisimulation-based representations in offline reinforcement learning.
\newblock \emph{Advances in Neural Information Processing Systems}, 36:\penalty0 28311--28340, 2023.

\bibitem[Zisselman et~al.(2023)Zisselman, Lavie, Soudry, and Tamar]{zisselman2023explore}
Ev~Zisselman, Itai Lavie, Daniel Soudry, and Aviv Tamar.
\newblock Explore to generalize in zero-shot rl.
\newblock \emph{Advances in Neural Information Processing Systems}, 36:\penalty0 63174--63196, 2023.

\end{thebibliography}


\end{document}
